\section*{The Mexican Criminal Conflict}

 The internal war in Mexico is a criminal conflict driven by territorial disputes over trafficking routes and land. Drug trafficking is not a new phenomenon but over the last decade it has affected all levels of Mexican society.  After the fall of the Colombian cartels in the 1990s, the landscape of violence related to drug trafficking completely shifted in Mexico as cartels gained new territorial control. Since this time, Mexican drug cartels have become the largest foreign supplier of methamphetamine and marijuana to the United States, effectively dominating the drug market. In fact, estimates claim that the drug trade employs over half a million people and generates roughly 4\% of Mexico's annual GDP.\footnote{\cite{shirk2011drug}} 

Although Mexican drug cartels have controlled the drug trade for decades, it was not until the 2006 election of Felipe Calder\'{o}n that drug-related violence began to soar and civilians found themselves under fire. In 2006, Calder\'{o}n became president and ushered in a new policy against the cartels. With support from the United States, the Mexican government initiated a massive campaign to combat drug-related violence.  Violence soared and between 2006 and 2011 and homicides nearly tripled from 10,452 to 27,213.\footnote{According to Mexico's National Statistics Institute (INEGI).} Sending armed actors into an already armed, violent, and competitive situation, Calderon's strategy became known as a failure. It did not address the fundamental needs of civilians or establish trusted local institutions where citizens could seek support in the realms of justice and security.  Instead, these policies complicated the security situation even more and created an unstable environment for reporters, government officials, and civilians.

The failure of ``Calder\'{o}n's War'' is partially attributable to the fact that DTOs are complex, with overlapping rivalries, family histories, splintered subgroups, and territorial disputes that drive their violent methods of political action.   DTOs are also engaged in extensive corruption networks across different levels of the government and throughout the Mexican territory. The influx of federal troops into areas of high criminal activity added further complication to pre-existing corruption. Because police in Mexico receive low pay (about \$9,000 to \$10,000 a year), their loyalty can often be bought by cartels; however, when bribery doesn't work, cartels routinely punish government officials with violence.\footnote{\cite{nathaniel2013}} Since combat and corruption between federal troops and cartels began, over forty mayors and numerous government officials have been murdered while increasing numbers of missing persons have been reported across Mexico as a result of cartels' increasing use of kidnapping.  Government corruption, civilian victimization, and a silenced media are severe problems deeply embedded in the conflict. 

Because of the cartels' brutal methods of punishment and gain, journalists and other forms of citizen representatives have been hesitant to report on these events. Journalists have not only been afraid to report out of fear that they \textit{might} be punished; in fact, they have been targeted and killed numerous times. In 2010, Carlos Santiago, an intern photographer for the Mexican newspaper El Diario, based in Ciudad Juarez, was shot and killed.  This was the second journalist from El Diario to be targeted.  The other was Armando Rodr�guez, a writer who worked the police beat and was killed in front of his own home. Following these deaths, the newspaper's editor drafted a plea to drug traffickers asking why they were being targeted. The article was published on the front page of the paper.\footnote{For a full interview with the editor see \cite{eldiario}.} Then, on April 28, 2012, Regina Mart\'{i}nez, a journalist for the national news outlet Proceso, was found dead in her home in Xalapa, Veracruz. This series of murders is indicative of a larger phenomenon across Mexico. According to the International Press Institute and the Mexican journalists' group ``Periodistas de a Pie,'' 103 journalists have been killed between 2000-2015 and 25 have disappeared. Since 2010, Mexico has been considered as deadly for journalists as Iraq; yet, these crimes continue with impunity. The Mexican case thus presents a relevant, timely, and difficult case for measuring the evolution of nuanced relationships between different violent actors. This study describes how the investigation of these relationships is possible. 

\section*{Data Challenges in The Mexican Case}
The quality of data on the Mexican criminal conflict remains mixed and generally suffers from underreporting. We know that there have been several key actors in this conflict over the years, including the Gulf Cartel, Juarez Cartel, La Familia Michoacana, Los Zetas, Sinaloa Cartel, and the Tijuana Cartel. However, because Mexican drug cartels are sometimes in conflict with one another and infiltrated by government officials, it is difficult to attribute responsibility for homicides or other violent events to one cartel or actor versus another. Although the noisiness of this data might seem daunting, it presents an opportunity for researchers to explore how they may improve data and knowledge about violent situations in contexts where it is often dangerous to do the costly on-the-ground ``legwork'' that is generally necessary to accrue such information. 

At present, the majority of data on violence in Mexico is based on homicide rates. Homicide data is produced from four main sources: Mexico's National Institute of Statistics and Geography (INEGI), the National System of Public Security (SNSP), the Mexican Federal Government, and La Reforma. In the beginning of the conflict (typically demarcated by Calder\'{o}n's assumption of the Presidency), national newspapers carried death counts related to drug violence.  \textit{La Reforma}  continues to maintain drug-related homicide data; however, transparency behind the methodology of this data collection remains uncertain. It is not known, for example, how the newspaper decides whether a homicide is drug-related or not. Mexico's INEGI has data based on death certificates, which allows one to acknowledge the manner of death (such as bullet wound). This data, however, is unable to attribute which homicides are linked to crime and which are unrelated. The National System of Public Security also has crime data based on local prosecutor reports, but its reliability is questionable due to the mixed incentives for governments to accurately report information. Finally, the federal government also has released data known as the ``Database of Alleged Homicides Related to Organized Crime.'' This database has information on executions and violence against authorities. Altogether, these data present several difficulties: first, they are not updated in real-time. To better understand the heterogenous evolution of civil conflict, researchers need to be able to describe conflict dynamics as they unfold. An additional, major criticism is that these data do not further our understanding about who is directly or indirectly responsible for these crimes.\footnote{This information summarizes an article will fuller details on the subject in \textit{Letras Libras}. See \cite{ley2012}. Melissa Dell uses this government data to assess whether or not PAN victories divert drug traffic to alternative routes predicted by the shortest paths in a networked trafficking model. The aim of \cite{dell2011}, however, is not to create actor-based event networks.} 

Acknowledging the shortfalls of pre-existing data, our analysis improves upon existing data by providing actor event data.  While we can only provide a rough estimation of actors involved in each conflictual event, this is a considerable advancement from the current status quo of knowing little to no information about which actors are engaged in which violent events in Mexico.\footnote{The only other data similar to this format is from a project of machine-coded data of Spanish newspapers created by Javier Osorio and Alejandro Reyes. This data is currently unavailable.} 


\subsection*{ICEWS Data and The Mexican Criminal Conflict}

A key goal for this study is to leverage machine-coded reports to construct a network of armed actors that represents conflict over time in Mexico. In order to construct this network study,  we use the ICEWS actor-coded event data. The ICEWS event data is part of a larger project designed to operate as a crisis warning system for policymakers.\footnote{For a summary, see \cite{OBrien2010}.}	This database has enabled policymakers and researchers to forecast conflictual events around the world.\footnote{\cite{ward2013}} The machine-coded event data are gleaned from natural language processing of a continuously updated harvest of news stories, primarily taken from Factiva$^{TM}$, an open source archive of news stories from over 200 sources around the world. The baseline event coder is called JABARI, a java variant of TABARI (Text Analysis By Augmented Replacement Instructions) which has been developed by Philip Schrodt and colleagues.\footnote{(see \url{http://eventdata.psu.edu/})} This approach combines a ``shallow parsing'' technology of prior coders with a richer exploitation of syntactic structure.\footnote{This has increased accuracy (precision) from 50\% to over 70\%, as demonstrated in a series of ongoing (informal) evaluations of its output by human graders. Peak human coding performance is reported to be around 80\% \citep{king2003automated}.}

The models create each data point by obtaining three components of the news story: the sender of the event (i.e., who initiated the action), the receiver or target of this action, and then the event type itself. We subsetted this data according to relevant ``violent'' cameo codes in order to gain access to all events relating to any armed actors such as rebels, insurgents, government, and the police. These events, in essence, capture any type of violent conflict between different actors. The event type itself is coded according to the Conflict and Mediation Event Observation (CAMEO) ontology.\footnote{See \cite{gerner2009} for the full summary of the project.} The main distinguishing feature of CAMEO is its use of mediation related event codes. CAMEO does not assume that a meeting is a peaceful interaction, for example, but is able to decipher whether meetings between actors are related to mediation, or negotiation. CAMEO also includes four categories for violence (structural violence, unconventional violence, conventional force, and massive unconventional force) as well as a rich system of sub-categories.

To begin to understand how to leverage the ICEWS data for country-level network analysis, we have taken a subset of data from the larger ICEWS corpus.  We constructed a SQL query to gain data subsetted according to all four ``violent" cameo codes as well as any actions related to all armed actors such as rebels, insurgents, government, and the police. Through the process of reviewing and cleaning the ICEWS data in preparation for my analysis, we encountered two key problems with the data. The first problem relates to the vague nature of the actor names in the data, which we improve upon via manual re-coding. The second problem we identify incentivizes ICEWS programmers to improve the parsing algorithm used in the creation of the original data.  While there are many unique actor names, the bulk of these descriptions are likely too vague for network construction. For example, the majority of cases relating to criminal violence use descriptors such as: ``Armed Gang,"  ``Armed Opposition,"  ``Attacker," ``Hitman," ``Drug Gang,'' ``Armed Band,'' and ``Criminal." A number of other cases have actor names such as ``Men'' or ``Citizens" as well as Military descriptors. 


To improve on these issues, we subset the data to include the raw text available from the larger ICEWS  database. Then, using a subset of cases from 2004-2013, we reviewed each individual case. This task includes two main goals: to label events as they relate to specific drug cartel actors in the area and to address aggregation problems in the data. In addition, we coded for duplicate cases.  We identified that this data has fewer duplication issues than previously found in other ICEWS data (such as protests) but has a number of complex aggregation and parsing problems. 

The end result is a collection of 1,052 actor-coded violent events from 2004-2012. The general trend of this data fits the pattern we observe in other data based on homicide rates. 

	%remake graphic?
	%discuss directionality (decreases # of  events) 
\newpage
\begin{figure}[h!]
\centering
\includegraphics[scale=.4]{Figure1Final}
\caption{Counts of conflictual events over time in Mexico 2004-2012.}
\label{fig:counts}
\end{figure}

 \newpage
\section*{Creating a Conflict Network in Mexico}
	%quarterly level
	%consolidated government actors, or say to discuss the nets we look at the dissaggreated nets
Next, we demonstrate how networks can capture the behavior of armed groups within the Mexican conflict. These networks, discussed below, reveal important relationships between actors, relationships that are otherwise missed by traditional dyadic frameworks. In a simple, descriptive and systematic way, we are able to show which DTOs were involved in the most conflict relevant to one another. We show that the Sinaloa cartel engages in violence against the highest number of competitors as well as against the federal government's armed forces. After walking through these connections and the ways in which they evolve over time, we turn to relationships  \textit{unobserved} in the raw data: i.e reciprocity and clustering.  

To begin, we create sociomatrices for each year of the cleaned data.  These sociomatrices can be thought of as a summary of interactions between all actors involved in conflictual events within a year. Given that there are $n$ actors in a year we construct an $n$ x $n$ sociomatrix $Y$. The number of conflictual dyadic interactions for any actor $i$ and $j$ is simply the number of events between those two actors during each given year. The resulting matrix is an undirected, symmetric matrix, as represented below.
 
\[
\left[
\begin{array}{cccccc}
& actor_{i} & actor_{j} & actor_{k} & actor_{l} & actor_{m}\\
 actor_{i} & 0  & 2  & 0 & 0 & 0 \\
actor_{j} &2  & 0  & 0 & 2 & 1 \\
actor_{k} & 0  & 0  & 0 & 0 & 4 \\
actor_{l} & 0  & 2  &0   & 0& 4 \\
actor_{m} & 0 & 1 & 4  & 4 &  0 \\

\end{array}
\right]
\]

\noindent These matrices reveal a variety of interesting dynamics in the data, as shown in the network graphs of Figure~\ref{fig:2005}, Figure~\ref{fig:2007}, and Figure~\ref{fig:2009}. First, Figure~\ref{fig:2005} shows the network graph for the first year of the data. Beginning as early as 2005 to 2007, there is a substantial increase in the number of actors in the network. For simplicity, the actors are simply coded as DTO or Government actors. Government actors include different branches of the armed forces and government.

\begin{figure}[h!]
\includegraphics[scale=.4]{20052006}
\caption{The evolution of the Mexican criminal conflict, 2005. Orange nodes are government actors, teal corresponds to criminal organizations. The links (grey lines) are weighted by the number of conflictual events for that year. }
\label{fig:2005}
\end{figure}


In 2006 and 2007, shown in Figure~\ref{fig:2007}, the number of conflictual events increase. Both federal agencies and municipal forces are more active in the conflict.  We also observe a jump in the number of cartels involved in the network, as new actors such as La Familia Micho\'acana, the Mexican Mafia, and the military begin to take part in the conflict. In the period between 2007-2009, the networks reflect a change in the government policy at that time. By 2006 Calder\'{o}n was elected president and implemented a militarized strategy, sending federal police and  troops into the most contested regions. By 2007, we see that the relationships become more complex as cartels begin to interact with one another. This sufficiently reflects what we now know to be true: following Calder\'{o}n's 2006 efforts, violence increased due to grueling competition between newly fragmented cartels. 


\begin{figure}[h!]
\includegraphics[scale=.4]{20072008}
\caption{The evolution of the Mexican criminal conflict, 2006-2007. Orange nodes are government actors, teal corresponds to criminal organizations. The links (grey lines) are weighted by the number of conflictual events for that year. }
\label{fig:2007}
\end{figure}

 In 2008, similar dynamics are in motion, but it is difficult to decipher the various roles played by the different levels of government. Earlier, both federal and municipal governments seem similarly involved, but by 2008 and 2009 the municipal government has a broader connectivity to diverse actors while the federal police seems to become more limitedly engaged with the Sinaloa Cartel. Overall, density increases across these networks overtime.  

To further assess how different levels of government change their role in the network over time, Figure~\ref{fig:eigen} presents the eigenvalue centrality for each actor. Eigenvector centrality is calculated by assessing how well connected an actor is to the other actors of the network. Specifically, those with the highest eigenvector centrality are well-connected to other actors who are also well-connected across the network.  In this case we can consider an actor with high eigenvector centrality to be a key player in the conflict.  


\begin{figure}[h!]

\includegraphics[scale=.4]{20092011}
\caption{The evolution of the Mexican criminal conflict, 2009-2011. Orange nodes are government actors, teal corresponds to criminal organizations. The links (grey lines) are weighted by the number of conflictual events for that year. }
\label{fig:2009}
\end{figure}


Figure \ref{fig:eigen} reveals an interesting story of coordination between different levels of government. In 2004 both the municipal police and the federal police have similar levels of eigen centrality. In 2005 they both decrease, with a larger decrease for the municipal police. In 2006 we see that both actors have very high centrality, indicating a deep involvement in the conflict. This reflects the Calder\'{o}n strategy: sending federal forces into local areas to help bolster security enforcement and combat the cartels. We then see these two actors diverge. Federal police stay highly central in the network in 2007 while the municipal forces seem to back down. At this year these two actors experience the biggest gap in centrality, signaling that coordination between the two levels of governments decreases as federal forces take over combat operations relative to municipal security forces. 

\begin{figure}[h!]
\raggedright
\includegraphics[scale=.4]{EigenCent}
\caption{Eigen centrality at the yearly level (2004-2009) for both municipal police and federal police }
\label{fig:eigen}
\end{figure}

By 2008, we see a stark shift as municipal forces are once again involved in the conflict while federal forces seem to back down. This is particularly interesting because at this time it becomes well known to the public and to the government that Calder\'{o}n's strategy has largely failed, causing more violence rather than stemming it. Figure~\ref{fig:eigen} shows that after federal forces are very active from 2006-2007, in 2008 they then recede and the municipal police become more active in the network. Finally, the year 2009 shows closer centrality scores than 2008 and 2007, signaling that perhaps these two levels of government begin to cooperate again. 

Without the re-structuring of the event data into networks we would not be able to see the ways in which different levels of governments interact over time. This approach offers a more complete picture of these dynamics than has previously been possible and enables future research to examine what factors influence government coordination within the network.

\section*{\textcolor{red}{SM's modeling section}}

\textcolor{red}{STUFF HERE ON THE ACTUAL MODEL. Why we chose it, whats cool about it, how no one has used it for intrastate conflict. How it can handle counts.}


\subsection*{Key Variables}
The dependent variable in our study is the conflict event network. At present, we include four predictor variables.

NEW IVs HERE!

%The first two are eigenvector centrality and betweeness scores. These variables are generated on the previous iteration of the network and thus act as a nodal time-varying covariate for each year in the model. The primary motivation here is to test whether actor-level effects, such as being a large ``sender" of the conflict, can best predict future conflict. Third, we include a control variable for whether or not the actor in the network is a Drug Trafficking Organization. Last, we include a protest variable. 

Our protest data comes from a new data set, ``The Mexican Protest Against Crime Database" collected by Ley (2015). This data records the number of protests against violence in a municipality during each month for the time period of 2006-2014. Creating a protest variable for the network analysis is straightforward: we simply match municipalities across the protest data and the conflict event data. In doing so, we create a protest count for the number of protests in a DTO's region of operation. 

\begin{figure}[h!]
\centering
\includegraphics[scale=.4]{protestSandra}
\caption{Counts of protest events over time in Mexico 2006-2014.}
\label{fig:counts}
\end{figure}




